% Proposed canonical-manuscript excerpt. Labels and citation are namespaced.
% Include immediately after the existing 05-real-turning.tex section.
% Add bibliography-entry.tex to the manuscript's thebibliography environment.
% Full proofs and finite certificates are in the companion research report.
\section{The unique quartic transition and two radial turning points}
\label{rbif:int:section}
The remaining finite-coefficient question from the preceding section admits
an exact computer-assisted resolution. The following results are proved
in the companion report \cite{rbif:continuation}; the global assertion is
supported by an exhaustive interval cover, not by a numerical sweep.

Write $p_n=n^{-a}H_{n-1}^{(b)}$, $\mu=p_3/(2p_2)$, and
$\eta_{a,b}(\rho)=\mu+K\rho^2+Q\rho^4+R\rho^6+O(\rho^8)$.
Let $a=A(b)$ be the unique quadratic threshold $K(A(b),b)=0$ for
$0<b<1$, and put $q(b)=Q(A(b),b)$. Here $A$ denotes the threshold,
not the polynomial used below.

\begin{theorem}[Unique quartic transition]
\label{rbif:int:quartic}
There is exactly one $b_\dagger\in(0,1)$ with $q(b_\dagger)=0$.
One has $q<0$ below $b_\dagger$, $q>0$ above it, and
$q'(b_\dagger)>0$. With $a_\dagger=A(b_\dagger)$, certified
rational enclosures include
\[
 0.99749378987342042107460<b_\dagger
       <0.99749378987342042107461,
\]
\[
 1.00143018096149513812935<a_\dagger
       <1.00143018096149513812936.
\]
At this point,
\[
 -3.368744536\cdot10^{-11}<R_\dagger<-3.368744535\cdot10^{-11},
\]
\[
 -8.315139844\cdot10^{-10}
 <\det D(K,Q)(a_\dagger,b_\dagger)<-8.315139843\cdot10^{-10}.
\]
\end{theorem}
\begin{proof}[Proof outline with replay dependencies]
The exact interval verifier covers $[0,9/10]$ by $2048$ cells and
$[9/10,1]$ by $256$ cells. It verifies all $4609$ endpoint/midpoint
root brackets and implicit-graph barriers, and proves
\[
 q'>10^{-6}\quad(0\le b\le9/10),\qquad
 q''<-29\cdot10^{-6}\quad(9/10\le b\le1).
\]
It also certifies $q(9/10)<0<q(999/1000)$; $q(1)=0$ is exact.
Increasing behavior on the first interval and strict concavity on the
second therefore imply exactly one interior zero. A rationally
preconditioned contraction gives a unique simultaneous zero of $K,Q$ in
the report's narrower 40-decimal rectangle, and direct interval evaluation
gives the displayed nondegeneracy signs. The full inclusion arithmetic,
logarithm/exponential tails and the replay programs
\texttt{verify\_global.py}, \texttt{verify\_local.py} are in the report.
\end{proof}

For completeness the new coefficient is obtained without fitting.
Set
\[
 A_1(y)=4p_3y^2-4p_4y+p_5,\qquad
 B_1(y)=8p_4y^3-12p_5y^2+6p_6y-p_7,
\]
\[
 C_1(y)=16p_5y^4-32p_6y^3+24p_7y^2-8p_8y+p_9.
\]
Then the exact coefficient identity is
\[
 R=-\frac{A_1'(\mu)Q+4p_3K^2+B_1'(\mu)K+C_1(\mu)}{2p_2}.
\]
At $K=Q=0$, this reduces to $R=-C_1(\mu)/(2p_2)$.

\begin{theorem}[Local two-extremum region]
\label{rbif:int:fold}
The map $(a,b)\mapsto(\kappa,\lambda)=(K,Q)$ is an analytic local
coordinate system at $(a_\dagger,b_\dagger)$. There is an analytic curve
\[
 \kappa=\varphi(\lambda),\qquad
 \varphi(\lambda)=\frac{\lambda^2}{3R_\dagger}+O(\lambda^3),
\]
and a fixed sufficiently small radius interval such that the region
$\lambda>0$, $\varphi(\lambda)<\kappa<0$ has exactly two critical
points in that interval: a strict minimum followed by a strict maximum.
At fixed $b=b_\dagger$, as $a\uparrow a_\dagger$ from below,
\[
 \rho_{\max}(a,b_\dagger)
 =\left(\frac{K_a(a_\dagger,b_\dagger)}{3R_\dagger}\right)^{1/4}
    (a_\dagger-a)^{1/4}+O((a_\dagger-a)^{3/4}).
\]
\end{theorem}
\begin{proof}
Put $t=\rho^2$, $\eta=\Phi(t)$ and $U=\Phi_t$. Then
$U=\kappa+2\lambda t+3R(\kappa,\lambda)t^2+O(t^3)$.
Since $U_{tt}=6R_\dagger<0$ at the base point, $U$ is strictly
concave locally. Its unique local maximum crosses zero on the analytic
curve displayed above. Between that curve and $\kappa=0$, two simple
positive roots occur, with the stated extremum types. Rescaling
$t=(a_\dagger-a)^{1/2}s$ at fixed $b=b_\dagger$ proves the
fourth-root law. The report proves the complete local phase classification
and separately certifies an explicit rational pair with at least two
extrema below $\rho=1/\sqrt{4000}$.
\end{proof}

\begin{corollary}[Constant-curve rigidity]
\label{rbif:int:rigidity}
For $a\ge1$, $b>0$, the normalized-radius curve is constant on a
nonempty interval adjacent to zero if and only if $(a,b)=(1,1)$.
In that case $\eta_{1,1}=1/2$ on $0<\rho\le1$.
\end{corollary}
\begin{proof}
Constancy forces $K=Q=R=0$. The established quadratic classification
reduces $K=0$ to $(1,1)$ or the fractional threshold. On that threshold,
Theorem~\ref{rbif:int:quartic} makes $Q=0$ possible only at the unique
transition, where $R<0$. Finally
$F_{1,1}(z)=\tfrac12\log^2(1-z)$ gives the stated constant curve.
\end{proof}

\paragraph{Scope.}
Quartic uniqueness is global in $b$ along the quadratic threshold.
The two-extremum count is local in radius and parameters. The separate
rational witness certifies at least two extrema, not a global count on the
unit disk. Neither the full-radius integer-outer-order conjecture nor
constant-curve classification for $0<a<1$ is settled here. The report also
proves all-order reversion and M\"obius-transport identities, including the
first order-nine symmetry defect at $(a_\dagger,b_\dagger)$.
